\uuid{wV4b}
\exo7id{3780}
\titre{exo7 3780}
\auteur{quercia}
\organisation{exo7}
\datecreate{2010-03-11}
\isIndication{false}
\isCorrection{false}
\chapitre{Espace euclidien, espace normé}
\sousChapitre{Endomorphismes auto-adjoints}
\cpge{mpsi-mp2i.calcul-matriciel-systemes.e, pcsi.determinants.c, ptsi.determinants.a, mp-mpi.reduction-endomorphismes.b, pc.reduction.a, psi.reduction-endomorphismes.a, pt.reduction-endomorphismes.a, mp-mpi.reduction-endomorphismes.g, pc.complements-espaces-vectoriels.d, psi.complements-espaces-vectoriels.d, bcpst2.produit-scalaire.c}
\module{Algèbre}
\niveau{L2}
\difficulte{}

\contenu{
\question{
Soit $A \in \mathcal{M}_n(\R)$ symétrique telle qu'il existe $k \in \N^*$ tel que
$A^k = I$. Montrer que $A^2 = I$.
}
}
